% ============================================================================
% Chapter 27 --- Wave 9: Lifecycle automation and maturity
% ----------------------------------------------------------------------------
% Purpose : self-service provisioning, GitOps discipline at scale, mature
%           operational stage. Last wave of the playbook.
% Page target: ~4 pages.
% Dependencies: Waves 1-8 in some form (D17.3); reference layer Ch. 14
%               (orchestration plane four properties D14.1, composition
%               discipline D14.4, no admissible third path D14.7).
% Mirrors wave template (D18.2).
% ----------------------------------------------------------------------------
% Decisions registered in _coherence-EN.md:
%   D27.1 self-service is rolled out per consumer group, not institution-wide
%   D27.2 GitOps repository becomes the operational source of truth
%   D27.3 ticket-driven legacy flows preserved during transition
%   D27.4 exit criteria reference Orchestration conformance subset + service
%         metrics
% ============================================================================

\chapter{Wave~9 --- Lifecycle automation and maturity}
\label{ch:wave9-lifecycle}

\begin{caveatbox}[Dependencies]
\noindent Wave~9 applies the wave chapter template established in
\trchap{ch:wave0-foundations}{}. The required pre-requisite is
the closure of Waves~1--8 in some form: not every wave must have
reached its full scope, but each must have closed the elements
on which Wave~9 builds. It applies the orchestration contract of
\trchap{ch:orchestration}{} (the four properties of D14.1, the
composition discipline of D14.4, the no-third-path rule of
D14.7). Closure is gated by the four Orchestration tests of
\trchap{ch:conformance}{}.
\end{caveatbox}

Wave~9 is the wave in which the architecture acquires the
operational ergonomics of a mature platform. The provisioning
that earlier waves have established as technically possible
becomes a self-service capability for the institution's
consumers. The GitOps discipline that earlier waves have
respected becomes the operational source of truth. The wave is
consequently both the last wave of the playbook and the wave in
which the institution stops thinking about ``the playbook'' and
starts thinking about routine operations.

%-----------------------------------------------------------------------------
\section{Goal and dependencies}
\label{sec:w9-goal}

The goal of Wave~9 is to make the architecture's day-to-day
operation sustainable without continuous engineering intervention.
Three concrete objectives follow:

\begin{itemize}[leftmargin=1.5em,nosep]
  \item self-service provisioning of TACs by their authorised
    consumers (researchers requesting research workspaces,
    teaching staff requesting lab environments,
    administrative staff requesting compliance workspaces),
    with the policy plane gating the request;
  \item GitOps repository elevated to the operational source of
    truth: changes to the architecture's running state are made
    by repository commit, reviewed under the same discipline as
    application code, and reconciled into running state by the
    orchestration plane;
  \item operational stability metrics agreed and reported, so
    that the institution can recognise when the architecture is
    operating at maturity.
\end{itemize}

%-----------------------------------------------------------------------------
\section{Pre-requisites}
\label{sec:w9-prereq}

The closure of Waves~1--8 is the principal pre-requisite. Three
operational conditions strengthen the wave when present:

\begin{itemize}[leftmargin=1.5em,nosep]
  \item the orchestration plane's GitOps repository is in active
    use by the engineering team, with established review and
    merge cadences;
  \item the policy plane has accumulated sufficient policy
    coverage that consumer-initiated provisioning requests do
    not require ad hoc human authorisation in the routine case;
  \item the conformance suite of \trchap{ch:conformance}{} has
    been exercised at least twice, so that the institution has
    a baseline for the maturity metrics that Wave~9 will use.
\end{itemize}

%-----------------------------------------------------------------------------
\section{Coexistence pattern}
\label{sec:w9-coexistence}

Wave~9's coexistence is between ticket-driven provisioning
(the institution's incumbent service-management pattern) and
self-service provisioning (the wave's target pattern).

\paragraph{Per-consumer-group rollout.} Self-service is rolled
out group by group rather than institution-wide. Typical
sequence: a research group with mature DevOps practice first,
then teaching environments where the requester is the
instructor, then administrative services where the requester is
typically a unit manager. Each group's rollout is evaluated
against the maturity metrics before the next group begins.

\paragraph{Ticket-driven legacy preserved.} Throughout the wave,
the existing ticket-driven flows remain available. Consumers
who do not yet have self-service for their use case continue
through tickets; the engineering team operates both flows
without requiring consumers to choose. The ticket-driven flow
is gradually retired only when the corresponding self-service
flow has demonstrated stability for an institutionally agreed
period (indicatively three to six months).

%-----------------------------------------------------------------------------
\section{Trajectory: greenfield}
\label{sec:w9-greenfield}

For institutions in a greenfield position, Wave~9 establishes
self-service and GitOps from the start of the institution's
production operation. There is no incumbent ticket-driven flow
to preserve, and the wave is correspondingly compressed.
Indicative duration is four to six months for the technical
configuration; the maturity metrics require additional time to
accumulate.

%-----------------------------------------------------------------------------
\section{Trajectory: brownfield single-suite-centric}
\label{sec:w9-brownfield-ms}

For institutions whose existing service-management is built
around ServiceNow, Microsoft Service Management, or a
comparable ITSM platform, Wave~9 introduces the self-service
capability alongside the ITSM platform rather than in
replacement of it.

\paragraph{Pattern.} The ITSM platform continues to operate for
the use cases that genuinely benefit from its workflow features
(approval chains involving multiple stakeholders, financial
authorisation, regulated processes). Self-service is introduced
for the use cases where the policy plane's authorisation is
sufficient and where the platform's workflow features add
ceremony rather than value. The ITSM platform may itself become
a consumer of the policy plane and the orchestration plane: an
ITSM-initiated request becomes a GitOps commit through a thin
adapter.

\paragraph{What changes in Wave~9.} The boundary between
self-service and ITSM-driven provisioning is documented and
governed; the institution's consumers know which path applies
to their request. The GitOps repository becomes the
authoritative state, with the ITSM platform acting as a
trigger rather than as the system of record.

\paragraph{Indicative duration.} Twelve to twenty-four months,
dominated by the per-consumer-group rollout and by the
ITSM--GitOps boundary work.

%-----------------------------------------------------------------------------
\section{Trajectory: brownfield hybrid}
\label{sec:w9-brownfield-hybrid}

For institutions with a multi-vendor compute estate and an
organisational-unit-level service-management heterogeneity, Wave~9
operationalises a unified self-service interface across the
heterogeneous substrate.

\paragraph{Pattern.} A common self-service portal (or, where
preferred, a CLI plus IaC pattern) is exposed across the
substrates that Waves~3, 8 and 9 have brought under common
governance. The portal does not eliminate organisational-unit-level
autonomy in service management; it provides the institutional
baseline that organisational-unit-specific extensions build on.

\paragraph{What changes in Wave~9.} The institution acquires a
shared provisioning surface that respects organisational-unit
autonomy while delivering consistent classification, audit, and
sovereignty properties across the estate.

\paragraph{Indicative duration.} Twelve to twenty-four months.

%-----------------------------------------------------------------------------
\section{Reversibility criteria}
\label{sec:w9-reversibility}

Wave~9 is reversible at the per-consumer-group level for as
long as the ticket-driven legacy flow remains operational.

\begin{itemize}[leftmargin=1.5em,nosep]
  \item Per-group rollback: a consumer group whose self-service
    rollout produces unworkable behaviour is returned to the
    ticket-driven flow while the divergence is diagnosed.
  \item Self-service portal disable switch: a tested mechanism
    for switching the portal to read-only or fully off within
    minutes is part of the wave's deliverables.
  \item Manual fallback procedure: for the case in which both
    self-service and ITSM are unavailable, a documented manual
    procedure ensures that the architecture's critical
    operations can still be exercised.
\end{itemize}

%-----------------------------------------------------------------------------
\section{Exit criteria}
\label{sec:w9-exit}

Wave~9 closes when the four Orchestration tests of the
conformance suite of \trchap{ch:conformance}{}
(\S\ref{sec:cf-suite}) pass, when the agreed self-service
ratio is achieved, and when the institution's operational
stability metrics are within agreed thresholds.

\begin{itemize}[leftmargin=1.5em,nosep]
  \item Orchestration conformance subset (idempotency,
    deactivation reverses consequential changes, GitOps repo
    suffices to rebuild state, workflow portability across
    engines).
  \item Self-service ratio: the institutionally agreed
    proportion of provisioning requests is fulfilled through
    self-service rather than through tickets. The ratio is
    typically expressed per consumer group rather than as a
    single institutional figure.
  \item Operational stability metrics: agreed thresholds for
    incident frequency, change-failure rate, and deactivation
    completeness are met over an institutionally agreed
    window.
\end{itemize}

\paragraph{Beyond Wave~9.} The closure of Wave~9 marks the end
of the playbook and the beginning of routine operations. The
sovereignty grid (\trchap{ch:sovereignty-grid}{}) continues to
apply at every renewal; the conformance suite
(\trchap{ch:conformance}{}) is exercised on its established
cadence; the international peer comparison anticipated in
phase~2 begins to operate against the institution's
longitudinal record.

\begin{realworldbox}[Self-service patterns]
\noindent Self-service provisioning patterns at research-
intensive institutions are reported in NREN community fora and
in EDUCAUSE publications. Some institutions operate
sophisticated self-service portals for research-cloud
provisioning (typically those participating in NSF ACCESS or
EGI federations); others operate hybrid arrangements with
ITSM-mediated approval for sensitive provisioning. The
specific portal designs and approval thresholds vary
substantially with the institution's research mix and
governance culture; no single design is universally
appropriate.
\end{realworldbox}

%-----------------------------------------------------------------------------
\section{Regulatory anchoring}
\label{sec:w9-regulatory}

\noindent Wave~9 closes the supply-chain and lifecycle loop. It
is the principal anchor of the supply-chain security obligation
of the applicable national NIS2 transposition (Art.~21(2)(d) of
the directive) operationalised through SLSA / Sigstore attestations
and through provenance-aware orchestration. Under \loc{ai-regulation} (in the EU, Art.~53 of the AI Act ---
obligations for providers of general-purpose AI models, including
model documentation and provenance), the wave's lifecycle
discipline applies to AI artefacts. The wave operationalises 27001~A.8.30 (outsourced
development) and contributes to NIST CSF \textsf{GOVERN}
(GV.SC --- supply-chain risk management) and \textsf{IDENTIFY}
(ID.AM --- asset management). It also operationalises
COBIT~APO10 (managed vendors) and BAI03 (managed solutions
identification and build). Detailed mapping in
\trchap{ch:regulatory-mapping}{} \S\ref{sec:rm-nis2},
\S\ref{sec:rm-ai-act}, \S\ref{sec:rm-iso},
\S\ref{sec:rm-nist-csf}, \S\ref{sec:rm-cobit}.
